\documentclass[10pt,a4paper]{amsart}
\usepackage[margin=19mm]{geometry}
\usepackage{amsmath,amssymb,mathtools}
\usepackage[scr=rsfs,cal=euler]{mathalfa}
\usepackage{xfrac}
\usepackage[unicode,hidelinks,bookmarks=false]{hyperref}
\newcommand{\ie}{i.e.\ }
\newcommand{\cf}{cf.\ }
\newcommand{\resp}{respectively\ }
\newcommand{\Subsection}{\subsection}
\providecommand{\nobreakdash}{\nobreak\hbox{-}}
\setlength{\emergencystretch}{2em}
\multlinegap0pt
\usepackage{amssymb}
\usepackage{mathtools}
\usepackage[shortlabels]{enumitem}
\newtheorem{lemma}{Lemma}
\newtheorem{proposition}{Proposition}
\newtheorem{remark}{Remark}
\newtheorem{corollary}{Corollary}
\title{Projector additive group codes}
\author{Javier de la Cruz}
\date{Source: arXiv:2604.15158v3 (CC BY 4.0)}
\begin{document}
\maketitle

\noindent\emph{Source and licence.} Projector additive group codes. Version arXiv:2604.15158v3 (CC BY 4.0), distributed by arXiv under the Creative Commons Attribution 4.0 International licence, http://creativecommons.org/licenses/by/4.0/, as recorded in the arXiv licence metadata for this exact version and shown on the abstract page https://arxiv.org/abs/2604.15158v3. Full licence notice: \texttt{licenses/arxiv-2604.15158-CC-BY-4.0.txt}.\par\smallskip

\noindent\textit{Editorial scope.} Counted unit: the article's corollary cor:idempotent-neu (source0.tex lines 2215--2223). The article's complete original proof of it is delivered with it (source0.tex lines 2225--2326, one \texttt{proof} environment of 102 lines). No mathematics is written, repaired or re-derived: the statement and the proof are the article's own lines, in the article's own notation, and no retained line is substituted or re-flowed.  Retained uncounted as the source-local context the proof needs: lines 927--936; lines 1004--1019; lines 1050--1068; lines 2064--2070; lines 2142--2148.  Nothing was omitted from any retained span, so the package carries no omission record.  No retained line contains a citation, so the wrapper carries no bibliography and no bibliography entry.  No retained line contains a figure environment, an image-inclusion command or drawing code, so no image is delivered and the package declares no asset dependency.  Editorial formatting only, in addition: an AMS \texttt{amsart} preamble that restates the article's own declarations and macros used by the retained lines, namely the environments \texttt{lemma}, \texttt{proposition}, \texttt{remark}, \texttt{corollary} on separate counters, together with the standard packages those lines need.  The retained environments print in this document as Lemma 1, Proposition 1, Remark 1, Corollary 1 in order of appearance, whereas the article prints the same items with the numbering of its own full document.  The complete native source of the article as distributed in the arXiv e-print for this exact version is bundled unchanged as \texttt{sources/198584/source0.tex}.
\begin{lemma}\label{lem:adjoint-basic}
Let $\langle\cdot,\cdot\rangle$ be a non-degenerate $F$-bilinear form on $KG$, and let $T,S\in \operatorname{End}_{FG}(KG)$. Then the following hold:
\begin{enumerate}
    \item $(T^\ast)^\ast=T$.
    \item $(TS)^\ast=S^\ast T^\ast$.
    \item $(\operatorname{Im}T)^\perp=\ker(T^\ast)$.
    \item $(\ker T)^\perp=\operatorname{Im}(T^\ast)$.
\end{enumerate}
Here $T^\ast$ denotes the adjoint of $T$ with respect to $\langle\cdot,\cdot\rangle$.
\end{lemma}

\begin{proposition}\label{idempotent-projector}
Let $e\in KG$ satisfy $e^2=e$. Then the following hold:
\begin{enumerate}
    \item Right multiplication by $e$ defines an $FG$-linear projector $\rho_e:KG\longrightarrow KG$ by $\rho_e(x)=xe$.
    \item $
\operatorname{Im}(\rho_e)=KGe$ and $
\ker(\rho_e)=KG(1-e).
$ Consequently,
$
KG=KGe\oplus KG(1-e)
$
as a direct sum of left $FG$-modules.
\item $
\rho_e(FG)=FGe$ is an additive left group code.
\end{enumerate}
\end{proposition}

\begin{remark}\label{propiedades-adjunto}
\begin{enumerate}
\item Let $e\in KG$ be an idempotent. Then $e^\ast$ is also an idempotent, since $(e^\ast)^2=(e^2)^\ast=e^\ast$. Consequently, right multiplication by $e^\ast$ defines an $FG$-linear projector $\rho_{e^\ast}:KG\longrightarrow KG$, given by $\rho_{e^\ast}(x)=xe^\ast$.

If $\rho_e^{\ast_{\mathrm E}}$ and $\rho_e^{\ast_{\mathrm{TE}}}$ denote the adjoints of $\rho_e$ with respect to $\langle\cdot,\cdot\rangle_{\mathrm E}$ and $\langle\cdot,\cdot\rangle_{\mathrm{TE}}$, respectively, then $\rho_e^{\ast_{\mathrm E}}=\rho_{e^\ast}$ and $\rho_e^{\ast_{\mathrm{TE}}}=\rho_{e^\ast}$. In particular, $\rho_e^{\ast_{\mathrm E}}=\rho_e^{\ast_{\mathrm{TE}}}$.

If $m$ is even and $\rho_e^{\ast_{\mathrm H}}$ and $\rho_e^{\ast_{\mathrm{TH}}}$ denote the adjoints of $\rho_e$ with respect to $\langle\cdot,\cdot\rangle_{\mathrm H}$ and $\langle\cdot,\cdot\rangle_{\mathrm{TH}}$, respectively, then $\rho_e^{\ast_{\mathrm H}}=\rho_{\overline e^\ast}$ and $\rho_e^{\ast_{\mathrm{TH}}}=\rho_{\overline e^\ast}$. In particular, $\rho_e^{\ast_{\mathrm H}}=\rho_e^{\ast_{\mathrm{TH}}}$.

\item Using Proposition~\ref{idempotent-projector} and  Lemma~\ref{lem:adjoint-basic}, a direct computation for the Euclidean and trace-Euclidean cases yields:
\begin{enumerate}
    \item $\operatorname{Im}(\rho_e) = \ker\big(\rho_e^{\ast_{\mathrm E}}\big)^{\perp_{\mathrm E}} = KG\,e.$
    \item $\ker(\rho_e) = \operatorname{Im}\big(\rho_e^{\ast_{\mathrm E}}\big)^{\perp_{\mathrm E}} = KG\,(1 - e).$
    \item $\operatorname{Im}\big(\rho_e^{\ast_{\mathrm E}}\big) = \ker(\rho_e)^{\perp_{\mathrm E}} = KG\,e^{\ast}.$
    \item $\ker\big(\rho_e^{\ast_{\mathrm E}}\big) = \operatorname{Im}(\rho_e)^{\perp_{\mathrm E}} = KG\,(1 - e^{\ast}).$
\end{enumerate}

Analogous identities hold for the Hermitian and trace-Hermitian forms.
\end{enumerate}
\end{remark}

\begin{lemma}\label{lem:adjoint-neu}
Let $\star\in\{\mathrm{TE},\mathrm{TH}\}$, where the case $\star=\mathrm{TH}$ is considered only when $m$ is even. Let $P,Q\in \operatorname{End}_{FG}(KG)$ be $FG$-linear projectors. Then the following are equivalent:
\begin{enumerate}
    \item $P\sim_{\mathrm{MvN}}Q$.
    \item $P^{\ast_\star}\sim_{\mathrm{MvN}}Q^{\ast_\star}$.
\end{enumerate}
\end{lemma}

\begin{proposition}\label{prop:equivalent-projectors}
Let $P,Q\in \operatorname{End}_{FG}(KG)$ be $FG$-linear projectors. Then the following are equivalent:
\begin{enumerate}
    \item $\operatorname{Im}(P)\cong \operatorname{Im}(Q)$ as left $FG$-modules.
    \item $P\sim_{\mathrm{MvN}}Q$ in $\operatorname{End}_{FG}(KG)$.
\end{enumerate}
\end{proposition}

\begin{corollary}\label{cor:idempotent-neu}
Let $e,f\in KG$ be idempotents, and let $\rho_e,\rho_f:KG\longrightarrow KG$ be defined by $\rho_e(x)=xe$ and $\rho_f(x)=xf$ for all $x\in KG$. Then the following are equivalent:
\begin{enumerate}
    \item $KGe\cong KGf$ as left $KG$-modules.
    \item $\rho_e\sim_{\mathrm{MvN}}\rho_f$ in $\operatorname{End}_{KG}(KG)$.
    \item $e\sim_{\mathrm{MvN}}f$ in $KG$.
    \item $KGe^\ast\cong KGf^\ast$ as left $KG$-modules.
\end{enumerate}
\end{corollary}

\begin{proof}
Since $\operatorname{Im}(\rho_e)=KGe$ and $\operatorname{Im}(\rho_f)=KGf$, the equivalence $(1)\Leftrightarrow(2)$ follows from Proposition~\ref{prop:equivalent-projectors}.

We next prove $(3)\Rightarrow(1)$. Assume that $e\sim_{\mathrm{MvN}}f$ in $KG$. Then there exist $u,v\in KG$ such that $e=uv$ and $f=vu$. Define maps $\varphi:KGe\to KGf$ and $\psi:KGf\to KGe$ by $\varphi(x)=xu$ for all $x\in KGe$ and $\psi(y)=yv$ for all $y\in KGf$. These maps are well defined. Indeed, if $x=ae\in KGe$, then
$
\varphi(x)=aeu=a(eu).
$
Since $e=uv$ and $f=vu$, we have
$
eu=(uv)u=u(vu)=uf,
$
and therefore
$$
\varphi(x)=a(uf)=(au)f\in KGf.
$$
Similarly, if $y=bf\in KGf$, then
$
\psi(y)=bfv=b(fv).
$
Since
$
fv=(vu)v=v(uv)=ve,
$
it follows that
$$
\psi(y)=b(ve)=(bv)e\in KGe.
$$ Moreover, suppose that $ae=a'e$ for some $a,a'\in KG$. Then  $aeu=a'eu$ i.e.
$\varphi(ae)=\varphi(a'e)$.  Similarly, if $bf=b'f$ for some $b,b'\in KG$, then $\psi(bf)=\psi(b'f)$.
Furthermore, both maps are left $KG$-linear. If $x=ae\in KGe$, then
$$
\psi(\varphi(x))=\psi(xu)=(xu)v=x(uv)=xe=(ae)e=ae=x.
$$
Likewise, if $y=bf\in KGf$, then
$$
\varphi(\psi(y))=\varphi(yv)=(yv)u=y(vu)=yf=(bf)f=bf=y.
$$
Thus $\varphi$ and $\psi$ are mutually inverse left $KG$-module isomorphisms, and therefore $KGe\cong KGf$ as left $KG$-modules.

We now prove $(1)\Rightarrow(3)$. Suppose that $\varphi:KGe\to KGf$ is an isomorphism of left $KG$-modules, and let $\psi:KGf\to KGe$ be its inverse. Set $u=\varphi(e)$ and $v=\psi(f)$. Then $u\in KGf$ and $v\in KGe$. Since $\varphi$ and $\psi$ are left $KG$-linear, for every $a,b\in KG$ we have
$
\varphi(ae)=au
$
and
$
\psi(bf)=bv.
$
Now $\varphi(e)=u$ implies
$
u=\varphi(e^2)=e\varphi(e)=eu,
$
and similarly
$
v=\psi(f^2)=f\psi(f)=fv.
$
Since $u\in KGf$, there exists $a\in KG$ such that $u=af$, and thus
$$
e=(\psi\circ\varphi)(e)=\psi(u)=\psi(af)=av.
$$
Using $u=af$, we obtain
$
uv=(af)v=a(fv)=av=e.
$
Similarly, since $v\in KGe$, there exists $b\in KG$ such that $v=be$, and hence
$$
f=(\varphi\circ\psi)(f)=\varphi(v)=\varphi(be)=bu.
$$
Using $v=be$, we obtain
$
vu=(be)u=b(eu)=bu=f.
$
Therefore $e\sim_{\mathrm{MvN}}f$ in $KG$.

Finally, we prove $(2)\Leftrightarrow(4)$ with respect to $\star=\mathrm{TE}$. By Lemma~\ref{lem:adjoint-neu}, we have
$
\rho_e\sim_{\mathrm{MvN}}\rho_f
$
if and only if
$
\rho_e^{\ast_{\mathrm{TE}}}\sim_{\mathrm{MvN}}\rho_f^{\ast_{\mathrm{TE}}}.
$
By Proposition~\ref{prop:equivalent-projectors}, this is equivalent to
$
\operatorname{Im}(\rho_e^{\ast_{\mathrm{TE}}})\cong \operatorname{Im}(\rho_f^{\ast_{\mathrm{TE}}})
$
as left $KG$-modules. Moreover, by Remark~\ref{propiedades-adjunto},
$
\rho_e^{\ast_{\mathrm{TE}}}=\rho_{e^\ast}
$
and
$
\rho_f^{\ast_{\mathrm{TE}}}=\rho_{f^\ast}.
$
Hence
$
\operatorname{Im}(\rho_e^{\ast_{\mathrm{TE}}})=KGe^\ast
$
and
$
\operatorname{Im}(\rho_f^{\ast_{\mathrm{TE}}})=KGf^\ast.
$
Therefore $(2)\Leftrightarrow(4)$.
\end{proof}

\end{document}
